\addtocontents{toc}{\protect\Needspace{3\baselineskip}}
\originalchapter{恒星结构、燃烧与演化}\label{ch:stellar}
\originalsection{静水平衡与能量预算}
恒星长期发光却没有立即坍缩, 说明压力梯度平衡重力, 同时有能量来源补充辐射损失. 球对称、慢于动力学时间变化的恒星以 $r$ 为半径, $m_r$ 为内部质量, $L_r$ 为穿过球面的光度. 质量守恒、力平衡及能量守恒分别为
\begin{equation}\label{eq:stellar}
\frac{\dd m_r}{\dd r}=4\pi r^2\rho,\qquad
\frac{\dd P}{\dd r}=-\frac{Gm_r\rho}{r^2},\qquad
\frac{\dd L_r}{\dd r}=4\pi r^2\rho\epsilon.
\end{equation}
$\epsilon$ 是单位质量净产热率, 演化模型中还须计中微子损失与热储能变化. 这三式尚不封闭: 需要状态方程、能量输运规律以及不透明度和核反应率. 中心条件为 $m_r,L_r\to0$, 表面则与大气边界匹配; 不能把给定总质量与任意中心温度当作互相独立的自由参数.

粗略地 $P_c\sim GM^2/R^4$、$\rho\sim M/R^3$. 理想气体 $P\sim\rho\kb T/(\mu m_p)$ 于是给 $T_c\sim GM\mu m_p/(\kb R)$. 这解释引力收缩为什么提高温度. 引力能 $W\sim-GM^2/R$, 热平衡下维里关系 $2K+W=0$ 使总能 $E=W/2<0$. 星体丢失能量后 $E$ 更负, $K$ 反而增大; 自引力系统具有这种负热容行为. 仅靠收缩供能的时间 $t_{\rm KH}\sim GM^2/(RL)$ 远短于太阳的核燃烧寿命.

\originalsection{辐射扩散与对流}
当平均自由程 $\ell=(\kappa_R\rho)^{-1}$ 远小于温度变化尺度, 各方向强度仅有小各向异性. 将 $B_\nu[T(s)]$ 按一个自由程展开并作角平均, $\int\cos^2\theta\dd\Omega=4\pi/3$, 得 $F=-c\nabla u/(3\kappa_R\rho)$. 使用 $u=a_RT^4$ 后,
\begin{equation}
F_r=-\frac{4a_RcT^3}{3\kappa_R\rho}\frac{\dd T}{\dd r},\qquad
\nabla_{\rm rad}\equiv\frac{\dd\ln T}{\dd\ln P}
=\frac{3\kappa_RL_rP}{16\pi a_RcGm_rT^4}.
\end{equation}
Rosseland 平均是按输运权重定义的不透明度调和平均, 不是任意频率不透明度的算术平均. 若气团上移时热交换慢, 它按绝热梯度 $\nabla_{\rm ad}$ 冷却. 环境在同一压强下温度下降更快, 即 $\nabla_{\rm rad}>\nabla_{\rm ad}$, 则气团较热、密度较低并继续上升. 这就是均匀组成下的 Schwarzschild 对流判据; 有平均分子量梯度时需加入组成稳定项.

\originalsection{核反应与演化尺度}
质子间库仑势垒使低能粒子难以接近, 量子隧穿使反应仍可发生. 非共振带电粒子截面可写为\[\sigma(E)=\frac{S(E)}E\exp\left[-\sqrt{\frac{E_G}E}\right].\] Maxwell 分布的 $\exp(-E/\kb T)$ 与隧穿因子竞争, 共同指数 $f(E)=E/\kb T+\sqrt{E_G/E}$ 在
$E_0=[E_G(\kb T)^2/4]^{1/3}$ 处最小. 对 $f$ 的二阶展开给出反应率主要来自 $E_0$ 附近的有限能段, 即 Gamow 窗口. 这里不将未知的 $S(E)$ 当作已推导常数, 详细核反应网络和实测截面不在本卷范围.

氢转为氦约释放反应质量的 $0.007c^2$. 若可燃质量占总质量的 $f$、氢质量分数为 $X$, 主序寿命估计 $t_{\rm nuc}\sim0.007fXM c^2/L$. 全对流低质量恒星可把更多氢带到燃烧区; 大质量恒星光度增长比质量快, 因而寿命短. 主序离开后的壳层燃烧、核心收缩和外层膨胀可在 H--R 图上区分, 但轨道的精确数值依赖完整输运、质量损失及组成演化, 不能用这一时间估计替代恒星演化程序.

\originalsection{平衡尺度与稳定性的区别}
估算恒星中心温度时, 必须区分维里总能与核燃料能. 对均匀球 $W=-3GM^2/(5R)$, 理想单原子气体总热动能 $K=3M\kb\bar T/(2\mu m_p)$. 维里关系给 $\bar T=GM\mu m_p/(5\kb R)$, 它是质量加权平均温度, 不等于中心温度. 中心温度通常更高, 要由具体密度与状态方程求出. 将此平均值的系数直接当作任意恒星中心公式会引入无依据的精度.

球的收缩和稳定也需要新的条件. 设均匀缩放时绝热 $P\propto\rho^\gamma$, 质量固定, 内能 $U\propto R^{-3(\gamma-1)}$, 引力能 $W\propto-R^{-1}$. 写 $E(R)=A R^{-q}-B/R$, $q=3(\gamma-1)$、$A,B>0$. 平衡 $E'(R)=0$ 给 $B=qA R^{1-q}$, 再代二阶导数得 $E''=q(q-1)A R^{-q-2}$. 因而这个同质缩放模型稳定要求 $\gamma>4/3$, 临界 $\gamma=4/3$ 对整体缩放边缘稳定. 这为相对论简并星质量极限提供物理解释, 但不是完整径向正常模分析, 对非均匀组成和强广义相对论星体须重新处理.

核燃烧还可给定温度反馈的直觉. 理想气体星体核心温度升高使反应率提高, 增加的热能推动膨胀, 密度与温度下降, 形成热调节. 简并核压力对温度不敏感, 热输入不能立即以相同方式膨胀调节, 可导致燃烧不稳定. 这解释氦闪的机制方向; 具体发生时刻和峰值必须使用相应反应率、热容和输运, 本卷没有编造未运行的氦闪曲线.
